\begin{frame}{``그래디언트 형태 동일''이 뜻하는 것}
  \begin{itemize}
    \item 마지막 줄(``형태 동일'')은 2.3절에서 \textbf{증명}할 내용.
      미리 눈여겨봐두자:
      \kb{시그모이드의 미분 $\sigma'(z) = \sigma(z)(1-\sigma(z))$ 가
      교차 엔트로피 손실의 미분과 만나서 통째로 약분}되어, 최종
      그래디언트가 선형회귀와 똑같은 $(h-y) \cdot \boldsymbol{x}$ 형태가 된다.
  \end{itemize}
  \vskip0.8em
  \begin{block}{``같은 도구를 다른 문제에 쓴다''의 구체적 뜻}
    \textbf{경사하강법의 코드 줄 자체를 그대로 복사해서,
    \texttt{h}를 계산하는 부분에 시그모이드 한 줄만 추가하는 것}.
    \\ 이 패턴 --- \textbf{모델 구조는 그대로 두고, 출력 해석과
    손실함수만 문제에 맞게 바꾼다} --- 는 학기 내내 반복.
  \end{block}
  \vskip0.4em
  {\footnotesize 2.3절 = 이 패턴의 첫 완전한 예(회귀$\to$분류).
  Week 9 신경망은 ``출력 변환''을 여러 층으로 늘린 것, Week 15
  생성모형은 같은 파이프라인을 ``생성''이라는 새 문제에 맞춘 것.}
\end{frame}
